\documentclass[exercice,a4]{thedocument}%
\usepackage[french]{babel}%
\usepackage{pstricks}%
\usepackage{pst-plot}%
\usepackage{multido}%
\startdatakeys
    \source{21GENMATAS1}
    \BO{2026}
    \cycle{4}
    \competencesDuSocle{CA,RA}
    \objectifsApprentissage{C4D1be,C4D2aa}
    \motsCles{vitesse,durée,lecture graphique,proportionnalité,fonction,fonction linéaire}
\enddatakeys
\begin{document}%

\medskip

Une mère et sa fille rentrent chez elles à pied en empruntant le même trajet de
10 kilomètres. La mère décide de s'y rendre en marchant et sa fille en courant.

Le graphique ci-dessous modélise les parcours de la mère et de la fille depuis leur départ.

\begin{figure}
\begin{center}
\psset{xunit=4cm,yunit=1cm,arrowsize=2pt 3}
\begin{pspicture}(-0.25,-1)(2.25,10)
\multido{\n=0.00+0.25}{10}{\psline[linewidth=0.15pt](\n,0)(\n,10)}
\multido{\n=0+1}{11}{\psline[linewidth=0.15pt](0,\n)(2.25,\n)}
\psaxes[linewidth=1pt]{->}(0,0)(0,0)(2.25,10)
\psline[linecolor=blue,linewidth=1.5pt](0,0)(0.25,3)(0.75,3)(1.833,10)
\psline[linecolor=red,linestyle=dotted,linewidth=1.5pt](0,0)(2,10)
\uput[r](0,10.2){Distance parcourue (en km)}
\rput(0.5,4.5){\blue Parcours de la fille}\psline[linecolor=blue]{->}(0.5,4.2)(0.3,3)
\rput(0.75,1.2){\red Parcours de la mère}\psline[linecolor=red]{->}(0.6,1.3)(0.45,2.2)
\uput[u](1.85,0){Temps de parcours (en h)}
\end{pspicture}
\end{center}
\end{figure}

\medskip

\begin{enumerate}
\item
	\begin{enumerate}
		\item Indiquer le temps mis par la mère pour rentrer chez elle, avec la précision que permet la lecture du graphique.
		\item Déterminer la vitesse moyenne en km/h de la mère sur l'ensemble de son parcours.
		\item La distance parcourue par la mère est-elle proportionnelle au temps ?
	\end{enumerate}
\item La fille est partie à $16$~h et est arrivée chez elle à $17$~h $50$. Elle a fait une pause durant sa course.
	\begin{enumerate}
		\item Indiquer la durée de la pause de la fille, avec la précision que permet la lecture graphique.
		\item Quand a-t-elle couru le plus vite: avant ou après sa pause ?
	\end{enumerate}
\item Combien de fois la mère et la fille se sont retrouvées au même endroit et au même moment, au cours de leur trajet ?
\item Dans cette question, on note $f$ la fonction qui, au temps de parcours $x$ (exprimé en heure) de la mère depuis le départ, associe la distance parcourue (exprimée en kilomètre) par la mère depuis le départ.

Parmi les propositions suivantes, recopier sans justification l'expression de $f(x)$ :

\[f(x) = \dfrac{1}{5}x \quad;\quad  f(x) = 5x \quad;\quad  f(x) = x + 5.\]

\end{enumerate}

\bigskip

\showthesource
\end{document}
